\section{Data Collection for Model Training}

We utilise 11 routing metrics for the feature importance analysis
ETX:
Hop count:
RSSI:

PPM: the total packet transmission per minute. These include
retransmission and packets from children nodes
CPU util: The percentage of active CPU over the total time
Drop rate: the number of successful transmission per one drop. The
higher the value the lower the drop ratio

Is new: which is the metric to differentiate between the current
parent or new parent, based on the way that we collect the metrics

In our evaluation, the nodes closer to root node would have higher
CPU util since they have to process more packets from children node,
thus also higher PPM and potentially higher Drop rate.

More over, we also include parent PPM, parent CPU util and parent
neighbour counts in the DIO to advertise to the child. Possibly to
link the load increase of the current node onto the selected parent

And finally choose 6 most importance metrics to train the model for
PDR prediction

The routing metrics are collected on two circumstances: on switching
parent and on receiving the DIO update from its current parent and
decide not to switch parent. Because of this, we add is new to the
training data feature to differentiate between current parent or new
parent as current parent also include the load of the current node.

Afterward, the training label PDR is calculated within a scope of a
node by calculating TX / RX between the current row and next data input time
the current row and the next row of the same node. We want the model
to be able to predict if this parent is selected, what would be the
potential PDR is.

Additionally, we filter the chunk having less than 30 seconds to
reduce the noise in
the training dataset.

The data is collected from simulation in COOJA with parameters include in Table ...

% TODO: add table for parameters here
